\section{泛函的延拓}
\label{sec:lfunc-extension}

一般的赋范空间没有内积，也不是积分空间，上一节的表示定理都用不上，我们甚至不知道$X^*$中是否有足够多的元素。在Hilbert空间中，任何非零向量$x$都被读数$\langle x,\cdot\rangle$识别出来，因为它在$x$处的值是$\lVert x\rVert^2\neq0$。一般空间中写不出这样的读数，但可以先在一个小子空间上定义它，再设法扩展到整个空间，并且不增大它的范数。本节先建立这种延拓，再用它证明连续线性读数能读出范数、区分向量。

\subsection{Hahn--Banach延拓}
\label{sec:lfunc-hahn-banach}

设$M$是$X$的线性子空间，$\ell$是$M$上的有界线性泛函，即$\lvert\ell(m)\rvert\leq\lVert\ell\rVert\lVert m\rVert$对$m\in M$成立。我们要找$L\in X^*$，使$L$在$M$上与$\ell$一致，并且$\lVert L\rVert=\lVert\ell\rVert$。延拓总能保持在$M$上的取值，难点在于不增大范数：$L$在$M$外的取值要受$\lvert L(x)\rvert\leq\lVert\ell\rVert\lVert x\rVert$的约束。

\begin{theorem}[Hahn--Banach定理]
\label{thm:lfunc-hahn-banach}
设$X$是赋范空间，$M\subset X$是线性子空间，$\ell:M\to\mathbb F$是有界线性泛函。则存在$L\in X^*$，使$L|_M=\ell$且$\lVert L\rVert=\lVert\ell\rVert$。
\end{theorem}

这一结论称为\term{Hahn--Banach定理}（Hahn--Banach theorem）。下面先完整证明延拓一维的核心步骤，再说明怎样推广到整个空间。

\paragraph{延拓一维}
设$\mathbb F=\mathbb R$，不妨设$\lVert\ell\rVert=1$（$\ell=0$时取$L=0$，否则除以$\lVert\ell\rVert$）。任取$z\notin M$，要把$\ell$延拓到$M\oplus\operatorname{span}\{z\}$上。线性延拓完全由一个数$t=L(z)$决定：$L(m+sz)=\ell(m)+st$。要求$\lVert L\rVert\leq1$，即
\[
\lvert\ell(m)+st\rvert\leq\lVert m+sz\rVert\qquad(m\in M,\ s\in\mathbb R).
\]
$s=0$时这就是$\lVert\ell\rVert\leq1$。$s\neq0$时两边除以$\lvert s\rvert$，并把$m/s$重新记为$m$，条件化为
\[
-\lVert m+z\rVert-\ell(m)\leq t\leq\lVert m+z\rVert-\ell(m)\qquad(m\in M).
\]

这样的$t$存在，当且仅当左端对$m$的上确界不超过右端对$m'$的下确界，即对一切$m,m'\in M$，
\[
-\lVert m+z\rVert-\ell(m)\leq\lVert m'+z\rVert-\ell(m').
\]
这确实成立：$\ell(m')-\ell(m)=\ell(m'-m)\leq\lVert m'-m\rVert\leq\lVert m'+z\rVert+\lVert m+z\rVert$，移项即得。所以可选的$t$组成一个非空闭区间，任取其中一个，就把$\ell$保范数地延拓了一维。

\paragraph{从一维到全空间}
若$X$可分，取可数稠密集$\{z_1,z_2,\ldots\}$，逐个把$z_k$添入定义域并重复上一步，得到$\operatorname{span}(M\cup\{z_k\}_k)$上范数仍为$\lVert\ell\rVert$的延拓。这个子空间稠密，而有界线性泛函一致连续，可以唯一地连续延拓到闭包即整个$X$上，范数不变：若$x_n\to x$，则$(L(x_n))$是Cauchy列，极限与所选序列无关，且$\lvert L(x)\rvert=\lim\lvert L(x_n)\rvert\leq\lVert\ell\rVert\lVert x\rVert$。不可分时，“逐个添加”要换成对所有保范数延拓按包含关系排序后使用Zorn引理，思路相同，这里不给出这一部分的证明。

\paragraph{复空间}
设$\mathbb F=\mathbb C$。复线性泛函由实部决定：记$u=\operatorname{Re}\ell$，则$\ell(m)=u(m)-\mathrm iu(\mathrm im)$，因为$\operatorname{Im}\ell(m)=-\operatorname{Re}(\mathrm i\ell(m))=-u(\mathrm im)$。把$X$看成实赋范空间，$u$是$M$上的实线性泛函，$\lvert u(m)\rvert\leq\lvert\ell(m)\rvert$，按实情形把$u$延拓为$U$且$\lVert U\rVert\leq\lVert\ell\rVert$，再令$L(x)=U(x)-\mathrm iU(\mathrm ix)$。直接验证$L$复线性且在$M$上等于$\ell$。对每个$x$取$\theta$使$e^{-\mathrm i\theta}L(x)=\lvert L(x)\rvert$，则$\lvert L(x)\rvert=L(e^{-\mathrm i\theta}x)$是实数，等于其实部$U(e^{-\mathrm i\theta}x)\leq\lVert\ell\rVert\lVert x\rVert$。所以$\lVert L\rVert\leq\lVert\ell\rVert$，反向不等式是显然的。

\paragraph{延拓可以不唯一}
一维延拓中的“非空区间”一般不止一点。在$(\mathbb R^2,\lVert\cdot\rVert_1)$中取$M$为横轴，$\ell(s,0)=s$，$\lVert\ell\rVert=1$。延拓写成$L_t(\bm x)=x_1+tx_2$，由第~\ref{sec:lfunc-dual-space}小节的对偶范数计算，$\lVert L_t\rVert=\max(1,\lvert t\rvert)$，所以每个$\lvert t\rvert\leq1$都给出保范数延拓。几何上，直线$x_1+tx_2=1$都经过$(1,0)$并支撑菱形单位球，菱形在这个顶点处有许多支撑线。换成欧氏范数，单位圆在$(1,0)$只有一条切线$x_1=1$，延拓唯一。定理保证的是存在，并且保持的只有两条约束：在$M$上的取值和范数上界，$M$外的取值一般不由它们确定。

\paragraph{区分向量与子空间}
延拓的一个直接用途，是用一个连续线性读数把一个向量与一个闭子空间分开。

\begin{theorem}[距离泛函]
\label{thm:lfunc-distance-functional}
设$M$是赋范空间$X$的闭子空间，$x_0\notin M$，$d=\inf_{m\in M}\lVert x_0-m\rVert$。则$d>0$，并且存在$L\in X^*$，使$\lVert L\rVert=1$，$L$在$M$上为零，$L(x_0)=d$。
\end{theorem}

\begin{proof}
$M$闭而$x_0\notin M$，所以$x_0$有一个邻域与$M$不交，$d>0$。在$M\oplus\operatorname{span}\{x_0\}$上令$\ell(m+sx_0)=sd$。$s\neq0$时$\lVert m+sx_0\rVert=\lvert s\rvert\lVert x_0+m/s\rVert\geq\lvert s\rvert d$，所以$\lvert\ell(m+sx_0)\rvert\leq\lVert m+sx_0\rVert$，$\lVert\ell\rVert\leq1$。取$m_n\in M$使$\lVert x_0-m_n\rVert\to d$，则$\ell(x_0-m_n)=d$，比值趋于$1$，故$\lVert\ell\rVert=1$。由定理~\ref{thm:lfunc-hahn-banach}延拓即得。
\end{proof}

\paragraph{\texorpdfstring{$\ell^\infty$}{l-infinity}上不来自\texorpdfstring{$\ell^1$}{l1}的泛函}
现在可以完成第~\ref{sec:lfunc-lp-dual}小节末尾的说明。记$c$为$\ell^\infty$中收敛数列组成的子空间，在其上令$\lambda(x)=\lim_nx_n$。由$\lvert\lim x_n\rvert\leq\sup\lvert x_n\rvert$及常数列，$\lVert\lambda\rVert=1$。由定理~\ref{thm:lfunc-hahn-banach}，$\lambda$延拓为某个$\Lambda\in(\ell^\infty)^*$。

若$\Lambda$来自某个$y\in\ell^1$，即$\Lambda(x)=\sum_ky_kx_k$，则对标准单位向量$\bm e_k$（只有第$k$项为$1$的数列）有$y_k=\Lambda(\bm e_k)=\lim_n(\bm e_k)_n=0$，于是$y=0$、$\Lambda=0$；可是$\Lambda$在常数列$(1,1,\ldots)$上的值为$1$，矛盾。所以$(\ell^\infty)^*$严格大于由$\ell^1$给出的那部分。这个$\Lambda$没有显式公式，它的存在完全依靠延拓定理；后面讨论弱收敛与反射性时还会再用到它。

\subsection{范数的对偶表示}
\label{sec:lfunc-norm-dual}

Hilbert空间中，$x\neq0$时单位读数$\langle x/\lVert x\rVert,\cdot\rangle$在$x$处恰好读出$\lVert x\rVert$。在一般赋范空间中，我们用延拓构造一个起同样作用的读数：先在$x$张成的一维子空间上规定它读出$\lVert x\rVert$，再保范数地延拓。

\begin{theorem}[范数的对偶表示]
\label{thm:lfunc-norming}
设$X$是赋范空间，$x\in X$，$x\neq0$。则存在$L\in X^*$，使$\lVert L\rVert=1$且$L(x)=\lVert x\rVert$。因此对一切$x\in X$，
\begin{equation}
\lVert x\rVert=\max_{\lVert L\rVert\leq1}\lvert L(x)\rvert ,
\label{eq:lfunc-norm-dual}
\end{equation}
并且若$L(x)=L(y)$对一切$L\in X^*$成立，则$x=y$。
\end{theorem}

\begin{proof}
在$M=\operatorname{span}\{x\}$上令$\ell(\alpha x)=\alpha\lVert x\rVert$，则$\lvert\ell(\alpha x)\rvert=\lVert\alpha x\rVert$，$\lVert\ell\rVert=1$。由定理~\ref{thm:lfunc-hahn-banach}延拓得到所需的$L$。另一方面，对一切$\lVert L'\rVert\leq1$有$\lvert L'(x)\rvert\leq\lVert x\rVert$，而上面的$L$取到等号，所以式~\eqref{eq:lfunc-norm-dual}中的上确界可以取到，是最大值；$x=0$时两边为零。最后，若$x\neq y$，对$x-y$取这样的$L$，得$L(x)-L(y)=\lVert x-y\rVert\neq0$。
\end{proof}

式~\eqref{eq:lfunc-norm-dual}说，一个向量的长度可以由全部单位读数的最大值恢复，向量本身也由全部读数唯一确定。这里的范数与对偶范数是对称的：泛函范数是单位球上读数的上确界，向量范数是单位泛函读数的最大值。

\paragraph{具体空间中的取范读数}
在具体空间中，取到最大值的读数常常可以直接写出。在$C([0,1])$中，若$\lvert f\rvert$在$t_0$处取到最大值，则$L=\overline{\operatorname{sgn}f(t_0)}\,\delta_{t_0}$满足$\lVert L\rVert=1$、$L(f)=\lVert f\rVert_\infty$，这里$\operatorname{sgn}z=z/\lvert z\rvert$。在实值$L^p$（$1<p<\infty$）中，对$f\neq0$取权重$g=\lvert f\rvert^{p-1}\operatorname{sgn}f/\lVert f\rVert_p^{p-1}$，第~\ref{sec:lfunc-lp-dual}小节的计算给出$\lVert g\rVert_q=1$、$L_g(f)=\lVert f\rVert_p$。定理的价值在于，即使写不出这样的读数，它也保证存在。

凸集之间的分离、优化问题的对偶和最优性条件，都是同一延拓思想的几何形式，属于优化问题，可参考\BookChapterRef{chap:constrained-optimization}{卷五第5章“约束优化与对偶”}。
